\documentclass{BeamerTemplate}

\title{Module 10 - Tests du khi-carré}

\subtitle{STT-1920 Méthodes statistiques}
\date{\today}

\begin{document}
\begin{frame}[t]
  \titlepage
\end{frame}

%%%%%%%%%%%%%%%%%%%%%%%%%%%%%%%%%%%
\section{Test d'ajustement}
%%%%%%%%%%%%%%%%%%%%%%%%%%%%%%%%%%%

\begin{frame}[t]{Test d'ajustement ($\chi^2$ goodness-of-fit)}

\begin{Boite}{Objectif}
Tester si une distribution observée est compatible avec une distribution théorique supposée.
\end{Boite}

\medskip
$H_0$ : les données suivent la distribution $F_0$ spécifiée.

\bigskip
Pour $k$ catégories avec effectifs observés $O_j$ et effectifs théoriques $E_j = n\,p_{0j}$ :

\begin{Boite}{Statistique de test}
$$\chi^2_0 = \sum_{j=1}^{k} \frac{(O_j - E_j)^2}{E_j} \;\dot{\sim}\; \chi^2_{k-1} \text{ sous } H_0$$
\end{Boite}

\medskip
Rejeter $H_0$ si $\chi^2_0 > \chi^2_{\alpha;\,k-1}$ ou valeur-$p < \alpha$.

\end{frame}

\begin{frame}[t]{Conditions d'application}

Pour que l'approximation $\chi^2$ soit valide :
\begin{itemize}\itemsep5mm
  \item Les effectifs \alert{théoriques} $E_j \ge 5$ pour chaque catégorie.
  \item Si ce n'est pas le cas, regrouper des catégories.
  \item $n$ suffisamment grand.
\end{itemize}

\bigskip
\underline{Exemple :} Dé équilibré — on lance un dé 60 fois. Est-il équilibré?
\begin{itemize}
  \item $H_0$ : $p_j = 1/6$ pour $j = 1, \ldots, 6$.
  \item $E_j = 60/6 = 10$ pour chaque face.
  \item Calculer $\chi^2_0 = \sum_{j=1}^{6}(O_j-10)^2/10 \sim \chi^2_5$.
\end{itemize}

\end{frame}

%%%%%%%%%%%%%%%%%%%%%%%%%%%%%%%%%%%
\section{Test d'indépendance}
%%%%%%%%%%%%%%%%%%%%%%%%%%%%%%%%%%%

\begin{frame}[t]{Test d'indépendance}

\begin{Boite}{Objectif}
Tester si deux variables qualitatives sont \alert{indépendantes} dans une population.
\end{Boite}

\medskip
Les données sont présentées dans un \alert{tableau de contingence} $r \times c$ :

\begin{center}
\begin{tabular}{c|cccc|c}
 & $B_1$ & $B_2$ & $\cdots$ & $B_c$ & Total \\ \hline
$A_1$ & $O_{11}$ & $O_{12}$ & & $O_{1c}$ & $n_{1\cdot}$ \\
$\vdots$ & & & & & \\
$A_r$ & $O_{r1}$ & & & $O_{rc}$ & $n_{r\cdot}$ \\ \hline
Total & $n_{\cdot 1}$ & & & $n_{\cdot c}$ & $n$ \\
\end{tabular}
\end{center}

$H_0$ : $A$ et $B$ sont indépendantes.

\end{frame}

\begin{frame}[t]{Effectifs théoriques et statistique de test}

Sous $H_0$ (indépendance) :
$$E_{ij} = \frac{n_{i\cdot} \times n_{\cdot j}}{n}$$

\begin{Boite}{Statistique du khi-carré d'indépendance}
$$\chi^2_0 = \sum_{i=1}^{r}\sum_{j=1}^{c} \frac{(O_{ij} - E_{ij})^2}{E_{ij}} \;\dot{\sim}\; \chi^2_{(r-1)(c-1)} \text{ sous } H_0$$
\end{Boite}

\medskip
Degrés de liberté : $\nu = (r-1)(c-1)$.

\medskip
Rejeter $H_0$ si $\chi^2_0 > \chi^2_{\alpha;\,(r-1)(c-1)}$ ou valeur-$p < \alpha$.

\end{frame}

\begin{frame}[t]{Exemple — tableau $3 \times 2$}

\underline{Exemple :} On veut savoir si le groupe sanguin et le sexe sont indépendants.

Tableau $3 \times 2$ ($r=3$, $c=2$) :
$$\nu = (r-1)(c-1) = (3-1)(2-1) = \alert{2}$$

\bigskip
\begin{itemize}\itemsep5mm
  \item Calculer les $E_{ij}$ à partir des totaux marginaux.
  \item Calculer $\chi^2_0 = \sum (O_{ij}-E_{ij})^2/E_{ij}$.
  \item Comparer à $\chi^2_{\alpha;\,2}$ (ex. $\chi^2_{0{,}05;\,2} = 5{,}991$).
\end{itemize}

\end{frame}

\begin{frame}[t]{Récapitulatif — tests du khi-carré}

\begin{center}
\begin{tabular}{lll}\hline
\textbf{Test} & \textbf{$H_0$} & \textbf{ddl} \\ \hline
Ajustement & Distribution $= F_0$ & $k-1$ \\[4pt]
Indépendance & $A \perp B$ & $(r-1)(c-1)$ \\ \hline
\end{tabular}
\end{center}

\bigskip
\begin{itemize}\itemsep5mm
  \item Les deux tests utilisent la même statistique $\chi^2_0 = \sum(O-E)^2/E$.
  \item Toujours vérifier que $E_{ij} \ge 5$.
  \item Test unilatéral droit seulement (on rejette quand $\chi^2_0$ est \alert{grand}).
\end{itemize}

\end{frame}

\end{document}
